% SGA 3, Expose VII_B, English reconstruction -- Loop 1.
% Scope: Section 5.1.1 in full on authority local pages 83--84;
% stops immediately before Section 5.1.2.
% Combined-reader pages: 596--597. Re-edition printed pages: 585--586.
% Original-edition marginal locator in scope: 552.
% Editorial note in scope: 145.
% Authority: Polo--Gille Exp7B-13oct24.pdf,
% SHA-256 D303690E0195654BBE7C3DAA5D89EA23E24CD709A3575217234AC19AE8836694.
% Jacob Reinhold's CC BY 4.0 Markdown was consulted as comparison lineage.
% No genuine diagram occurs in this component; all mathematics is native LaTeX.

\subsubsection*{5.1.1}
\phantomsection
\label{sga3:VIIB:5.1.1}

考虑如下范畴 \(\mathcal C\)。\(\mathcal C\) 的对象是由预有限 \(k\)-代数 \(A\) 和闭理想 \(J\subset A\) 组成的对 \((A,J)\)。\(\mathcal C\) 中的态射
\[
  \psi:(A,J)\longrightarrow(A',J')
\]
是将 \(J\) 映入 \(J'\) 的连续 \(k\)-代数同态
\[
  A\longrightarrow A'
\]

对应
\[
  (A,J)
  \longmapsto
  \bigl(\operatorname{Spf}(A/J),\operatorname{Spf}(A)\bigr)
\]
给出 \(\mathcal C\) 到由形式 \(k\)-簇 \(Y\) 和形式子簇 \(Z\subset Y\) 组成的对 \((Z,Y)\) 的范畴的反等价。态射
\[
  \varphi:(Z,Y)\longrightarrow(Z',Y')
\]
是将 \(Z\) 映入 \(Z'\) 的形式 \(k\)-簇态射
\[
  Y\longrightarrow Y'
\]

\(\mathcal C\) 的对象 \((A,J)\) 上的余群结构，是指 \(\operatorname{Spf}(A)\) 上满足下列条件的形式群结构\textup{（采用 \SGARefLink{sga3:VIIB:2.1}{2.1} 的记号）}：
\[
\begin{aligned}
  \textup{(1)}\quad&
  \Delta_A(J)
  \subset
  J\widehat\otimes_kA
  +
  A\widehat\otimes_kJ,
  \\
  \textup{(2)}\quad&
  \varepsilon_A(J)=0,
  \\
  \textup{(3)}\quad&
  c_A(J)\subset J.
\end{aligned}
\]
这些条件也就是说，
\[
  H=\operatorname{Spf}(A/J)
\]
是
\[
  G=\operatorname{Spf}(A).
\]
的形式子群。
\footnote{\label{sga3:VIIB:nde:145}若 \((A',J')\) 为 \(\mathcal C\) 的另一个余群，对应于形式 \(k\)-群对 \(H'\subset G'\)，则给定余群态射
\[
  (A',J')\longrightarrow(A,J)
\]
等价于给定将 \(H\) 映入 \(H'\) 的形式 \(k\)-群态射
\[
  G\longrightarrow G'
\]}

再假设 \(A\) 为局部的；等价地，\(\operatorname{Spf}(A)\) 为无穷小形式群。若 \(J\ne A\)，则条件 \textup{(2)} 和 \textup{(3)} 可由 \textup{(1)} 推出。事实上，\(A\) 的真闭理想 \(J\) 包含于增广理想 \(I_A\)，所以 \textup{(2)} 成立。此外，
\[
  M=\operatorname{Spf}(A/J)
\]
为 \(G\) 的形式子幺半群。由于 \(G\) 为无穷小的，\SGARefLink{sga3:VIIB:2.7:prop}{命题~2.7} 表明 \(M\) 为 \(G\) 的形式子群；因此 \textup{(3)} 也成立。
